\documentclass[11pt]{article}
\usepackage[T1]{fontenc}
\usepackage{mathptmx,amsmath,amssymb,amsthm,geometry,hyperref}
\geometry{margin=1in}
\hypersetup{colorlinks=true,linkcolor=blue,urlcolor=blue}
\newtheorem{theorem}{Theorem}
\newtheorem{corollary}[theorem]{Corollary}
\newtheorem{remark}[theorem]{Remark}
\title{Exact Ranking of Lucas Differences and Intersections of A191981--A191983}
\author{Qichao Wang}
\date{October 2026}
\begin{document}
\maketitle
\begin{abstract}
Let $L_0=2$, $L_1=1$, and $L_{n+2}=L_{n+1}+L_n$. We prove that every positive integer represented as $L_i-5L_j$ with $i,j\geq1$ has a unique such representation. We characterize the admissible index pairs, count the values generated by bounded indices exactly, and show that the number of positive values at most $x$ is $\frac12(\log_\varphi x)^2+O(\log x)$. We give an exact count using only three boundary blocks and quadratic formulas for the four nonzero residue subsequences. We construct two infinite families common to A191981 and A191982, and explain the initial zero.
\end{abstract}

\section*{The index region}
Write
\[
 S=\{L_i-5L_j>0:i,j\ge1\},\qquad
 A(x)=\#\{s\in S:s\le x\},\qquad
 \varphi=\frac{1+\sqrt5}{2}.
\]
The identities
\begin{equation}\label{eq:identities}
 L_{n+5}=5L_{n+1}+3L_n\quad(n\ge0),\qquad
 L_{j+3}-5L_j=-2L_{j-2}\quad(j\ge2)
\end{equation}
follow from the Lucas recurrence (or by checking two initial values and then applying that recurrence).

\begin{theorem}\label{thm:region}
For $i,j\ge1$, the difference $L_i-5L_j$ is positive if and only if
\[
 (i,j)=(4,1)\quad\text{or}\quad i\ge5\text{ and }1\le j\le i-4.
\]
For each $i\ge5$, these $i-4$ differences are distinct and lie in
\begin{equation}\label{eq:block}
 3L_{i-5}=L_i-5L_{i-4}\ \le\ L_i-5L_j\ \le\ L_i-5.
\end{equation}
\end{theorem}
\begin{proof}
The Lucas numbers increase strictly from index $1$ onward. For $j\ge2$, the second identity in \eqref{eq:identities} gives $L_{j+3}<5L_j$, whereas the first gives $L_{j+4}=5L_j+3L_{j-1}>5L_j$. Monotonicity settles all $i$ for such $j$. For $j=1$, one checks $L_3=4<5$ and $L_4=7>5$. Equation \eqref{eq:block} and distinctness within a fixed $i$ follow from monotonicity of $L_j$, $j\ge1$.
\end{proof}

\section*{Uniqueness and growth}
\begin{theorem}\label{thm:unique}
Every $s\in S$ has exactly one representation $s=L_i-5L_j$ with $i,j\ge1$.
\end{theorem}
\begin{proof}
The residues $(L_1,L_2,L_3,L_4)\pmod5$ are $(1,3,4,2)$ and repeat with period four; in particular, they are distinct. Suppose
$L_i-5L_j=L_k-5L_\ell>0$.
Reduction modulo $5$ forces $i\equiv k\pmod4$. If $i=k$, strict monotonicity of the Lucas numbers from index $1$ gives $j=\ell$.

Otherwise, by symmetry take $k\ge i+4$. Theorem~\ref{thm:region} gives $i\ge4$, $k\ge8$, and $\ell\le k-4$. Thus
\[
 L_k-5L_\ell\ge 3L_{k-5}\ge3L_{i-1}>L_i>L_i-5L_j,
\]
a contradiction. Here $L_n$ increases for $n\ge1$, and $3L_{i-1}>L_i=L_{i-1}+L_{i-2}$ for $i\ge4$.
\end{proof}

\begin{corollary}\label{cor:finite}
For every integer $N\ge4$, precisely
\[
 1+\sum_{i=5}^N(i-4)=1+\frac{(N-4)(N-3)}2
\]
positive values have representations with first index $i\le N$.
Furthermore, the elements of $S$ are nonzero modulo $5$, and within each fixed nonzero residue class modulo $5$, the blocks of values obtained from consecutive permissible first indices occur in strictly increasing order.
\end{corollary}
\begin{proof}
The count follows from Theorems~\ref{thm:region} and~\ref{thm:unique}. Nonzero residues follow from the four-term Lucas residue cycle. Two first indices in the same residue class differ by at least four. If $k\ge i+4$, the proof of Theorem~\ref{thm:unique} shows every value in the $k$-block exceeds every value in the $i$-block.
\end{proof}

\begin{theorem}\label{thm:asymptotic}
As $x\to\infty$,
\[
 A(x)=\frac12(\log_\varphi x)^2+O(\log x).
\]
Consequently, if $b_1<b_2<\cdots$ enumerates the positive elements of $S$, then
\[
 \log_\varphi b_n=\sqrt{2n}+O(1).
\]
\end{theorem}
\begin{proof}
Let $m(x)$ be the largest integer $m\ge1$ with $L_m\le x$ (for $x\ge3$). Binet's formula
$L_m=\varphi^m+(-\varphi)^{-m}$ implies
$m(x)=\log_\varphi x+O(1)$.
Every positive difference with $i\le m(x)$ is at most $x$, giving
\[
 A(x)\ \ge\ 1+\frac{(m(x)-4)(m(x)-3)}2
\]
for large $x$. Conversely, a positive difference with $i\ge5$ is at least $3L_{i-5}$ by \eqref{eq:block}, so $i\le m(x/3)+5$. This gives
\[
 A(x)\ \le\ 1+\frac{(m(x/3)+1)(m(x/3)+2)}2
\]
for large $x$. Since both $m(x)$ and $m(x/3)$ equal $\log_\varphi x+O(1)$, these bounds prove the asymptotic. Finally, substituting $x=b_n$ and using the uniqueness of all positive terms gives
$n=\frac12(\log_\varphi b_n)^2+O(\log b_n)$; solving for $\log_\varphi b_n$ gives the claimed estimate.
\end{proof}


\section*{Exact counting and residue-class unranking}
For integral $y$, write $H(y)=\#\{j\ge1:L_j\le y\}$ and $[z]_+=\max(z,0)$.
\begin{theorem}[Three-block exact counting]\label{thm:exact}
For every integer $x\ge11$, let $m$ be the unique integer satisfying
$L_m\le x<L_{m+1}$. Then
\begin{equation}\label{eq:exact}
\boxed{\displaystyle A(x)=1+\frac{(m-4)(m-3)}2+
\sum_{h=1}^{3}\left[m+h-4-
H\!\left(\left\lfloor\frac{L_{m+h}-x-1}{5}\right\rfloor\right)
\right]_+.}
\end{equation}
For $0\le x<11$, $A(x)$ counts the members of $\{2,3,6,9\}$ not exceeding $x$.
\end{theorem}
\begin{proof}
All admissible differences with first index $i\le m$ are smaller than
$L_i\le x$. For $i\ge m+4$, \eqref{eq:block} gives
\[
L_i-5L_j\ge3L_{i-5}\ge3L_{m-1}>L_{m+1}>x.
\]
Indeed, $3L_{m-1}-L_{m+1}=L_{m-1}-L_{m-2}>0$ because $m\ge5$.
The only remaining first indices are $i=m+1,m+2,m+3$. In one such block,
$L_i-5L_j\le x$ is equivalent to
$L_j>\lfloor(L_i-x-1)/5\rfloor$.
There are $i-4$ admissible indices $j$, so removing the ineligible ones gives
the positive-part expression in \eqref{eq:exact}.
\end{proof}

\begin{corollary}[Exact ranking and membership]\label{cor:algorithm}
The published OEIS A191981 has $a(1)=0$.
For $n\ge2$, its $n$-th term is the least integer $x$ for which
$A(x)\ge n-1$. A finite upper bound is $L_N$, with $N\ge4$ the least
integer such that $1+(N-4)(N-3)/2\ge n-1$.
For membership, $x>0$ determines the first index $i\bmod4$ by
$x\equiv L_i\bmod5$. The corresponding index blocks are disjoint
and ordered, so a single Lucas lookup for $(L_i-x)/5$ gives the
unique witness if one exists.
\end{corollary}
\begin{proof}
The first assertion follows from uniqueness and \eqref{eq:exact}.
At least $n-1$ positive differences have first index at most $N$ and
every such difference is smaller than $L_N$. For the last assertion,
the Lucas residues modulo $5$ distinguish $i\bmod4$.
If $k\ge i+4$, then
$\min B_k\ge3L_{k-5}\ge3L_{i-1}>L_i>\max B_i$,
so blocks with the same first-index residue occur in increasing order.
\end{proof}

\begin{theorem}[Quadratic residue-class enumeration]\label{thm:residue}
Let $S_r=\{s\in S:s\equiv r\pmod5\}$ for $r=1,2,3,4$,
each listed in increasing order. The blocks of $S_r$ have
first indices $i_r+4k$ ($k\ge0$), where
$(i_1,i_2,i_3,i_4)=(5,4,6,7)$. The cumulative numbers
of terms through block $k$ are
\[
\begin{array}{ll}
F_1(k)=(k+1)(2k+1),&F_2(k)=1+2k(k+1),\\
F_3(k)=2(k+1)^2,&F_4(k)=(k+1)(2k+3).
\end{array}
\]
Set $F_r(-1)=0$. To obtain the $n$-th term in $S_r$, choose
the least $k\ge0$ with $F_r(k)\ge n$, put
$t=n-F_r(k-1)$ and $i=i_r+4k$. Then
\[
s_{r,n}=L_i-5L_{i-3-t},
\]
except that $s_{2,1}=2$ (the isolated $i=4$ block).
\end{theorem}
\begin{proof}
The residue cycle determines $i_r$.
By Corollary~\ref{cor:algorithm}, successive blocks are ordered;
within a block the values increase as $j$ decreases.
The block sizes are respectively $1+4k$,
$4k$ for $k\ge1$ (with one exceptional value for $k=0$),
$2+4k$, and $3+4k$. Summing the lengths gives the four
quadratics. The $t$-th element in a nonexceptional block has
$j=i-3-t$. Inverting a quadratic inequality determines $k$ exactly.
\end{proof}

\section*{Relations to A191982 and A191983}
Let $T=\{5L_i-L_j>0:i,j\ge1\}$, the positive part of
A191982. Unlike $S$, it has repeated representations:
$2=5L_1-L_2=5L_3-L_6$ and
$4=5L_1-L_1=5L_2-L_5$.
\begin{theorem}[Two infinite intersection families]\label{thm:overlap}
The set $S\cap T$ is infinite. Specifically,
\begin{align*}
3L_n&=L_{n+5}-5L_{n+1}=5L_{n-1}-L_{n-5}
&& (n\ge6),\\
L_{n+6}-5L_n&=5L_{n+2}-L_{n-4}
&& (n\ge5)
\end{align*}
are two families of positive values common to A191981 and A191982,
and hence present in their union A191983.
\end{theorem}
\begin{proof}
Each displayed identity holds for the first two permitted indices
and is preserved by the Lucas recurrence, so it holds for all
the indicated $n$. Each value is positive by Theorem~\ref{thm:region}.
Finally, $3L_n$ is strictly increasing for $n\ge6$.
\end{proof}

\begin{remark}[The initial zero]
A191981 begins with $0$, but $5$ never divides a Lucas number: $L_i\equiv1,3,4,2\pmod5$ for $i\ge1$. Hence $L_i-5L_j=0$ has no solution. The Mathematica program linked from A191980 replaces nonpositive differences by zero before taking their union. Thus the published sequence is $\{0\}\cup S$, in increasing order; its zero is a program-generated sentinel, not a Lucas difference. In particular, the number of published terms at most $x$ is $A(x)+1$ for $x\ge0$.
\end{remark}

\paragraph{Reference.}
C. Kimberling, \emph{OEIS A191981: Ordered sequence of nonnegative differences $f-5g$, where $f$ and $g$ are Lucas numbers}, \url{https://oeis.org/A191981} (accessed October 2026). The related sequences are \url{https://oeis.org/A191982} and \url{https://oeis.org/A191983}; the Mathematica code appears under \url{https://oeis.org/A191980}.
\end{document}
